\section{Method}
\label{sec:methodology}

The method separates the conversion step from the measurements that follow it.
Conversion is the intervention; the evaluation then asks task-proximal questions
about proxy match, VLM grounding, material mechanisms, and stack use.

\subsection{Conversion as Intervention}
\label{sec:conversion_intervention}

The intervention stays within the same USD scene family. From a root USD file,
the tool discovers referenced layers, payloads, variants, and clips, then
writes sibling \texttt{\_noMDL} assets. Composition arcs stay intact and point
to those siblings; source files do not change.

Each processed file swaps MDL shaders for PreviewSurface. The rewrite covers
base color, roughness, metal, normal maps, and texture links. Existing paths
stay in their declaring layer; missing paths use fixed values. The rewrite
covers a practical subset of MDL; clearcoat, displacement, opacity, and
emission stay audit risks.

\subsection{Evaluation Checks}
\label{sec:evaluation_checks}

The first check screens four assets
from Isaac Sim with PSNR, SSIM, LPIPS, CLIP, \mbox{DINOv2}, and headless
load/render checks. It checks whether conversion keeps the main structure in one
controlled object case.

The second check is the VLM grounding protocol on GRScenes. The original GRScenes tree is
the source condition; noMDL counterparts are materialized in a scratch workspace
and paired with matched target-centered views and 2D target boxes. QA separates
clean PASS pairs from PASS/WARN pairs where the target is visible and materials
shift. Two local VLM probes share the schema: Gemma4 is primary;
\mbox{Qwen2.5-VL} is a second probe. We record label agreement,
normalized-1000 point hits, and raw pixel point hits. The raw pixel check marks
probe-specific coordinate use. The probes test fixed prompt-and-coordinate
behavior under two local backends.

The third check maps visible shifts to material bins with a selected baseline check. A
GRScenes row is used only when all three outputs are present: source MDL,
\mbox{ConvertAsset} noMDL, and NVIDIA PreviewSurface. Converted files must
contain zero active MDL, and the target must render. This check covers
selected bins; full semantic coverage remains open. Gaps stay as wrapper cases
for failures or known gaps.

The fourth check tests stack entry through the \mbox{KuJiaLe} route, which supplies
matched original/noMDL scenes for InternNav and
\mbox{DualVLN}~\cite{Wei2026Ground}; a separate check repeats official-scene
load/render trials. The result is limited to stack entry: converted scenes load
and run in the selected stack under the tested route.
Table~\ref{tab:acl_evidence_gate_registry} summarizes what each check uses and
what remains outside the experiment.

\begin{table*}[t]
\centering
\footnotesize
\setlength{\tabcolsep}{3pt}
\begin{tabular}{p{0.15\textwidth}p{0.30\textwidth}p{0.24\textwidth}p{0.22\textwidth}}
\toprule
Gate & Evidence used & Claim supported & Forbidden promotion \\
\midrule
Proxy similarity & Four Isaac Sim assets; 24 matched render pairs; PSNR, SSIM, LPIPS, CLIP, DINOv2, and load/render diagnostics. & Converted assets pass a visual and feature-screening gate. & Task preservation, broad robustness, or speedup. \\
VLM grounding & GRScenes 15-pair clean pilot and frozen 30-pair target-centered stress set with answer, point-in-box, and coordinate scoring. & Material perturbation can be evaluated as bounded answer and point-grounding evidence. & Population-level GRScenes benchmark claims or general VLM rankings. \\
Material mechanisms & Selected opacity/transparency, emission, normal/bump, and displacement/height bins, plus clearcoat and procedural supplemental cases; original/noMDL/NVIDIA conditions are separated. & Covered bins have selected static and qualitative support; clearcoat and procedural texture remain scoped failure or limitation cases. & NVIDIA population failure rates, official-scene NVIDIA performance, or procedural-texture success. \\
Embodied-data sanity & Official KuJiaLe InternNav 99 paired episodes over three scenes plus official-scene load/render stability. & Converted scenes can enter the selected official downstream stack and load/render repeatedly. & Broad embodied-navigation robustness, all-scene preservation, or noMDL speedup. \\
\bottomrule
\end{tabular}
\caption{Evidence-gate registry for the ACL/ARR claim boundary. Each gate lists the evidence that may be used, the claim it supports, and the promotion that remains forbidden without new evidence.}
\label{tab:acl_evidence_gate_registry}
\end{table*}

